
Folkebladet, June 15, 1898

The Free Church’s Second Annual Meeting.

—

(Continued from page 3.)

\subsection{Thursday Morning.}

Prayer meeting from 9:00 to 9:30, led by Pastor P. Winther.

Yesterday’s minutes were read and adopted.

The Credentials Committee reported through Pastor A. J. Løgeland that 52 pastors and professors, together with 80 voting delegates, had submitted their names.

A letter from Pastor S. Romsdahl, Calumet, Mich., was read.

Pastor A. Helland maintained his resignation as secretary.

Pastor H. A. Urseth was then elected secretary.

The meeting then took up for consideration the proposed “Rules for a Lutheran Free Church.” It was resolved to consider and vote upon the draft point by point.

“1. According to God’s Word, the congregation is the proper form of the Kingdom of God on earth.”

Prof. G. Sverdrup: It is of great importance that we express ourselves concerning these matters; for here the concern is not merely to adopt points, but that we also acknowledge from the heart what we confess in these Rules. It is something new when in Point 1 we declare the congregation to be the proper form of the Kingdom of God on earth; for through hundreds of years the Church has been called the proper form.

If we adopt this point, however, we go back to the time of the Apostles and to God’s Word.

According to God’s Word, the congregation is not a society founded by Pope, King, Council, or any other authority; rather, the congregation is the gathering of believing people around God’s Word and the means of grace. There is not a single indication in God’s Word that any governing body is to be placed over the congregation to rule over it.

Point 1 adopted.

“2. The congregation consists of believing people who, through the use of the means of grace and the gifts of grace according to the direction of God’s Word, seek their own salvation and eternal blessedness and that of all their fellow human beings.”

Adopted.

“3. According to the New Testament, the congregation requires an outward organization with a register of its members, election of officers, appointed times and places for its gatherings, etc.”

Pastor E. Hansen wished that the points might be discussed, so that more light could be shed upon the matter.

Pastor M. Midthun was of the same opinion; he especially wished the Committee to express itself.

Sverdrup maintained that it was a bad custom into which they had fallen, always waiting for the Committee to express itself.

As for the point, he would say that it was a very much disputed point among believers. Not all had come to the recognition that the congregation requires an outward organization, etc. Many believe that such an organization is harmful to the congregation. And this misunderstanding certainly has its basis in the lamentable conditions among our church people. It is a great misunderstanding when certain believers will not belong to the outward organization of the congregation, but leave the work of the congregation to worldly people—in order themselves to preach freely, as it is called.

The first Christians had a good accounting of the affairs of their congregations and did not leave them to worldly people. We even have recorded the names of the first members of the congregation (Acts 1:13–15)—a testimony for all times that believers were not found outside the congregation. The first Christians counted their people (Acts 2:41; 4:4). In Acts 6:2 it says that the Twelve called together the whole multitude of the disciples. They therefore knew who they were.

Many remain outside the congregation because they believe it is a dangerous organization. They feel more burden than support from the congregation; and therefore they withdraw. And there is reason enough to feel so and to speak ill of the congregation, such as it all too often appears among us. Nevertheless, it is necessary and biblical for believers to join themselves together for common edification and common work for the cause of the Kingdom of God. We must therefore hold fast to the congregation and ourselves stand within it.

Pastor Brøns: I am heartily in agreement with this paragraph. We must have a congregation, and a visible congregation, an outward organization. That there are people who are so spiritual that they will not belong to any outward congregation, we know; but we cannot work in a proper manner by standing outside. It becomes only something loose, something upon which one cannot really get a proper hold.

There is much that is visible about us human beings, and the congregation ought also to resemble us a little in this respect—not merely resemble God.

Pastor Thomasberg feared that it would not be so straightforward to realize these principles under our conditions; for not many congregations correspond to Point 2.

Pastor Midthun: Many questions arise when one is to begin organizing a congregation. Some will have believers only, and no others with them; but this view leads to spiritual conceit, which in turn leads to a craving to judge and rule, and finally results in complete dissolution. And then one has spiritual fanaticism. We must have a congregation according to God’s Word, with an organization in a definite place.

Pastor Hansen thought that the point distinguished between the congregation and the outward organization. He wished to have the words “for the sake of the work” added.

Point 3 adopted.

“4. The congregation’s outward organization commonly also includes people who in God’s Word are called dead and sleeping, for which reason within the congregation as well there is need of the counsel of awakening and the preaching of conversion, just as it is the holy calling of believers, also within the congregation, to be light and salt, so that outward connection with the congregation may not become a false comfort and a hindrance to true and living Christianity.”

Pastor Midthun: God’s Word would present a congregation without spot or wrinkle; nevertheless, we find that all the congregations of which we read had spot and wrinkle. It is not God’s fault, but the fault and infirmities of human beings. Our goal ought to be to attain the ideal which God’s Word sets before us. The speaker could not assent to the point as it stood, but wished to propose some changes.

Pastor Broen proposed that the point be omitted.

Jonas Hetland proposed that the word “hypocrites” be substituted for the word “dead.”

Pastor Aas supported the proposal that the point be omitted. In any case he wished it to be redrafted. His understanding was that it is the believers who are to gather themselves into a congregation, and he could not see that there is anything dissolving in this. The love which drives believers together cannot be dissolving or lead to dissolution and pride. The life of Christ was not such. It was humility. Where the Spirit and love of Christ are permitted to dwell, there is no dissolution or pride.

The first Christians, however, were not afraid to confess themselves believers. If the congregation is to become visible, this must take place by the believers coming together.

Prof. Sverdrup: We must take care that we do not come to go against ourselves and our own principles. We desire awakening; but if this point, which speaks of awakening, were to be omitted, it would give the appearance that after all we did not mean it, or did not dare to establish it when the matter was put to the test. If we desire a living, pure, believing congregation, then we must do something to preserve it living and pure.

We can easily, before we know it, enter upon the confessional principle that in the congregation all are believers, and all believers are in the congregation, and that for them the preaching of conversion is not needed. We must seek to labor in such a way that we do not work ourselves fast into such a position.

In the New Testament we find that the members of the congregation are continually admonished to keep themselves awake and not fall asleep. If it is to be established that a person is a Christian because he belongs to a congregation, or that a person is a Christian because he is a pastor, then we enter again upon the old unbiblical tracks.

The speaker was, moreover, also of the opinion that the point was unfortunately worded, and therefore proposed the following substitute:

“Since not all members of the congregation’s outward organization are always believing people, and since such hypocrites often take false comfort from their outward connection with the congregation, it is the holy calling of the congregation to purify itself both through the awakening proclamation of God’s Word, through earnest admonition and reminder, and through the exclusion of manifest and obstinate sinners.”

\subsection{Thursday Afternoon.}

Devotion by Pastor E. Broen.

The chairman then called the meeting to order. He then announced that the Augsburg Seminary Board of Trustees had received word that the Supreme Court of Minnesota had decided the Augsburg case in favor of Augsburg’s friends. This would surely bring joy to the hearts of the Augsburg people. That we had the moral right to the school none of us had doubted, and we could therefore now rejoice that truth and right had prevailed. Let us rejoice with trembling and thank God for his grace. Let the outcome become an inducement to us to support Augsburg Seminary, so that it may become a school for pastors in accordance with the requirements of the Free Church.

At the proposal of Pastor O. Paulson, the assembly rose and sang the hymn “Now Thank We All Our God.”

The chairman led in prayer.

The Credentials Committee then reported 20 new members as entitled to vote at the meeting. A number of advisory members were admitted, among them the brothers Holter, Pastors Oftby, Gaardsmoe, Winkler, and other members of Hauge’s Synod present and entitled to attend.

The Nominating Committee proposed the following committee on the Seminary Board’s report: Prof. Oftedal, Pastor E. Berlien, Pastor P. Winther, E. Nilsen, and Jonas Hetland. All elected.

For the standing Heathen Mission Committee: Prof. Sverdrup, Prof. Blegen, Pastor M. A. Pedersen, J. Berge, and Pastor E. Aas. All elected.

For the standing Inner Mission Committee the following were elected: Pastor Harbo, Pastor A. Helland, H. Rasmussen, J. H. Egeland, and Pastor E. M. Broen.

Ordination Committee for this meeting: Pastor Harbo, Pastor Midthun, Pastor Kjelaas, Pastor Løgeland, and Pastor P. Nilsen.

Standing Ordinator: Prof. Oftedal.

Organization Committee: Prof. Oftedal, Prof. W. Pettersen, and Pastor N. Halvorsen.

Auditing Committee: P. D. Olsen and O. A. Hain.

After these elections, discussion of Point 4 was again taken up.

Pastor Quanbeck: It is true, as has been said, that it is terrible that people should take the fact that they belong to a congregation as a certain sign that they are Christians, whether they are so or not. It is therefore not strange that many become apprehensive and remain outside the congregation. Here there are two extremes against which one must guard; but one must nevertheless not be so afraid of one extreme that one falls into another.

The speaker favored Sverdrup’s substitute, since it expressed his standpoint more clearly.

Pastor Arewik had always found a contradiction between Points 2 and 4. In any case, the expressions were misleading. He could not, however, support the proposal to put “hypocrites” in place of “dead.”

He had found it more difficult to convince hypocrites than openly ungodly people of the necessity of conversion. From the expressions “a whole multitude of dead” and “has left his first love,” he found that even in the apostolic congregation there were those who did not belong to the Lord. Nevertheless, the goal is that the congregation shall consist of believers.

Pastor Broen would vote for Sverdrup’s substitute. He would not have proposed that the point be omitted if he had known the substitute.

It had caused misunderstanding when in Point 2 one said that the congregation consists of believers, and then again in Point 4 that others also can be included.

We must express ourselves in such a way that we are not misunderstood.

He had, of course, not made his proposal because he was opposed to awakening, but because the former point was misleading.

Pastor Brønø: There is also a difference among hypocrites. Many find it easy to confess with the mouth, others more difficult. I am not prepared without more ado to acknowledge a person as a Christian merely because he smiles at me. Those who cannot so easily make confession are often more earnest Christians than certain facile talkers.

We can so easily be mistaken. I would therefore not take part in expelling from the congregation people who live honorable lives, even though I did not dare believe that they are living Christians. One has opportunity to purify the congregation through the proclamation of God’s Word and the preaching of conversion. Many believers are so carnal that when there is someone they do not like, they immediately want to get that person out of the congregation. We speak of the congregation as being without spot or wrinkle; but that is difficult to attain so long as sinful human beings are alive.

We have spoken much about the ungodliness of our congregations; but let each one examine himself, and we shall find that we are no better ourselves, but that we are among the most wretched.

Let there be more self-knowledge and humility.

Sverdrup’s substitute was adopted unanimously.

Points 5, 6, and 7 adopted. They read as follows:

“5. The congregation itself governs its own affairs under the authority of God’s Word and Spirit and acknowledges no other ecclesiastical authority or governing power over itself.

6. The free congregation esteems and loves all the spiritual gifts which the Lord bestows for its edification, and seeks to stir up the gifts and promote their use.

7. The free congregation receives
   [the source is defective here]
   the congregations can render one another in the work for the advancement of the Kingdom of God.”

“8. Such mutual assistance consists partly in this, that the spiritual powers of one congregation benefit another through free meetings, mutual visits, lay activity, etc.,

and partly in this, that the congregations voluntarily and according to the prompting of God’s Spirit cooperate for the accomplishment of such tasks as may exceed the powers of a single congregation.”

Pastor P. Nilsen: I had not thought to speak in this matter, but only to hear what the other brethren had to say; I would nevertheless ask leave to say a little in connection with this point. I have been in serious doubt whether we ought to adopt these principles at all; I feared that we might fasten ourselves to these paragraphs and come away from God’s Word.

In the course of my work in the congregations one thing has become clear to me, and that is that it is highly necessary that the simple truths of the Bible concerning the congregation be brought forward again.

All who proclaim the Word in the congregations must make earnest of this and help our people to greater light in this matter by again and again setting forth the truth of God’s Word concerning the congregation.

If the Free Church is to advance toward a sound development, it is necessary that there be free interaction among the congregations. But this has its particular difficulties. There has been so much factional activity among us that the fear of something similar among us stands blocking the way. It must be clearly understood that nothing of this kind is intended by this; rather, when a visit is made to a congregation, it is for the purpose of building up, not tearing down or splitting apart.

Many excuse themselves and prefer to avoid taking part in the work; they would rather see others, preferably pastors, do it. The fact that one meets with ill will also holds many back.

When the congregations find that the use of the gifts of grace brings blessing and not trouble, they will become more willing to use them.

These many discussion meetings which are now being held here and there will surely, little by little, prepare the way for the use of the gifts of grace. It is of great importance that these principles be realized.

Point 8 adopted.

Points 9, 10, and 11 adopted.

“9. Among such purposes are mentioned particularly a school for pastors, distribution of Bibles and other writings and periodicals, Inner Mission, Heathen Mission, Jewish Mission, deaconess homes, orphanages, and other works of mercy.

10. Free congregations have no right to demand that other congregations submit to their opinion, will, judgment, or determination, and therefore every dominion of a majority of congregations over a minority is to be rejected.

11. Such common organs as may be considered desirable for the cooperation of the congregations, such as larger or smaller meetings, committees, officers, etc., cannot in a Lutheran Free Church impose upon the individual congregation any duty, bond, compulsion, or burden, but are entitled only to come forward with proposals and appeals and persons.”

[The final wording of Point 11 appears incomplete or textually defective in the source.]

“12. Every free congregation has, like every individual believer, the prompting of God’s Spirit and love’s right to do good and to labor for the salvation of souls and the awakening of spiritual life, as far as its ability and strength extend; in such free spiritual activity it is restricted neither by parish bonds nor by denominational boundaries.”

Pastor Yberstab thought that this point might require a more exact explanation; he could not assent to the final part. If it is adopted, one has thereby established as a principle that one may without further ado intrude into other congregations and use it as an excuse for carnal motives. The speaker then cited a series of statements by Luther against intrusion into the congregations of others.

Pastor Telle proposed that the final part of the point be omitted.

Pastor M. A. Pedersen was exceedingly glad for the point; he did not want anything omitted. That such free and unrestricted activity would bring its difficulties is certainly true; but that cannot be helped. As for Luther’s statements, he would only say that if they went against the statements of Jesus, then it was best to hold to Jesus.

I understand it thus, that Jesus wills that we bear witness concerning him to all without exception. Wherever there is an unconverted person, it is the duty of God’s children to bring such a one to God, even though he should belong to another congregation than ours.

Pastor E. Hansen wished to ask where it stood in God’s Word that we are to let the unconverted go to hell because they belong to another congregation or another church body than ours.

Delegate Rørvik had the same question to ask.

Pastor Broen likewise wished to remind them that it does not stand in God’s Word that one is to let people go to hell if they belong to another church body. He nevertheless believed that one would attain the same result even if the final part of the point were omitted; it would then not be quite so provocative. One ought to seek to avoid offense. But he favored free activity. It might sound strange; but in this respect he was willing to be un-Lutheran.

Prof. Sverdrup: The point is both biblical and Lutheran. If someone sends out a missionary or an emissary, we surely will not call that wrong. The New Testament knows nothing of parish bonds or denominational boundaries.

We would therefore not act contrary to God’s Word if, for example, we sent someone to bear witness to our dear ones in the Fatherland, of whom we knew that they lived without fellowship with the Lord. That they might possibly shut him out would be another matter. But, for that matter, they do not do so in Norway. It is only here in America that one does such things.

Would it be wrong if we received into a hospital of ours a sick person from another church body? Let there be a free and open way. These parish bonds and denominational boundaries are something human beings have set up.

Luther, who has been cited here, stands before his Judge and not before our tribunal. If his statements are pleasing to God, then it is well; but we are not bound by everything Luther has said.

Moreover, we know from church history that Luther by no means kept so strictly to the old church order or within the old ecclesiastical boundaries. He was quite bold in that respect, yes, so bold that he was placed under the ban. If we become as bold as Luther, there will be no lack of boldness in the Free Church.

Point 12 was adopted in its entirety.

It was resolved that on Friday morning the draft of the congregational order should be taken up for consideration, and that a sufficient number of copies of the same should be printed for distribution among the members of the meeting.

The Credentials Committee then reported the following committee on the mission reports: Pastor O. Paulson, Pastor E. O. Larsen, Pastor H. Quanbeck, Mr. Fingal Fosliën, and Baard Andersen.

In the evening there was worship in Trinity Church by Pastor P. Nilsen, and an offering for the Inner Mission was received.

Pastor O. L. Torvik and Carl Amundsen in St. Luke’s Church, Pastor Risløv in St. Olaf’s Church, Pastor E. Hansen and M. Midtbøun in St. Peter’s Church.

(Continued.)
